\documentclass[11pt,a4paper]{article}
\usepackage[T1]{fontenc}
\usepackage[utf8]{inputenc}
\usepackage[english]{babel}
\usepackage{lmodern}
\usepackage[margin=27mm]{geometry}
\usepackage{microtype}
\usepackage[hidelinks]{hyperref}
\title{Detecting False Mission Completion from Robot Inspection Evidence: An Operator-Centered Study}
\author{Mateusz Pomianek}
\date{Version 0.3.0 -- 21 November 2026}
\begin{document}
\maketitle

\begin{abstract}
Operator decisions about mission completion can be tested using inspection reports that either conceal or explicitly expose missing and mismatched visual evidence. Does a structured evidence-and-omissions report reduce false acceptance of incomplete inspection missions compared with a success notification or binary status? The prospective design is as follows: Generate independently labelled complete, incomplete and wrong-object cases from versioned offline RVI fixtures. Randomize report presentation across declarative, binary and evidence-explicit formats with balanced information volume. The primary outcome is False-acceptance proportion among objectively incomplete missions, measured at the mission-decision level. The manuscript defines a research question and an evaluation strategy; it does not report experimental estimates.
\end{abstract}

\section{Introduction}
Operator decisions about mission completion can be tested using inspection reports that either conceal or explicitly expose missing and mismatched visual evidence. The scientific contribution under consideration is a falsifiable comparison in which alternative explanations, measurement limitations and distinct hypotheses from related ARC Robo studies can be assessed independently. Both the source of observations and the assumptions needed to interpret results must be documented before any confirmatory comparison.

\section{Related research and novelty assessment}
The following comparison anchors the scope in selected independently verifiable existing work. The citations are not evidence of completed experiments or of established originality:
\begin{itemize}
\item \textbf{Prior art 1:} Transparency improves the accuracy of automation use, but automation confidence information does not \cite{source1}. Shows transparency can influence automation decisions; the proposed RVI study must isolate image-evidence provenance and false completion rather than generic transparency or confidence.
\item \textbf{Prior art 2:} Trust as a design principle in human--robot collaboration \cite{source2}. Reviews HRI trust and explanation; it is not evidence that a specific inspection checklist prevents false mission acceptance.
\end{itemize}
\subsection{Specific test beyond existing work}
Marginal difference in false acceptance of objectively incomplete inspection missions between evidence-explicit and information-length-controlled status reports. This planned contrast is the candidate independent contribution. A literature overlap audit must confirm that prior publications have not already demonstrated an equivalent contrast under comparable measurement conditions.

\section{Primary estimand and prospective decision rules}
Marginal difference in false acceptance of objectively incomplete inspection missions between evidence-explicit and information-length-controlled status reports. Use participant/case crossed effects, a prespecified non-inferiority margin for true completion acceptance and an ethics-approved pilot to set precision and sample size. H1 is the primary scientific hypothesis unless the source protocol states a different frozen confirmatory contrast. H2 and H3 are supporting or exploratory until a prospective multiplicity strategy (e.g., Holm adjustment) is selected. Any sample size, quality boundary, non-inferiority margin, effect threshold or confidence limit not grounded in existing evidence is explicitly pending approval before confirmatory evaluation.

\section{Detailed scientific scope and novelty boundary}
The human endpoint is erroneous acceptance of an objectively unfinished robot inspection mission. The visual report can provide an unqualified completion message, binary status, or a provenance-linked inventory of required photographs and omissions. Mission truth is computed independently of what the interface shows. The study does not estimate objective robot inspection success or camera quality.

\section{Falsifiable hypotheses}
\paragraph{H1.} An explicit evidence-and-omissions report reduces false acceptance among incomplete missions compared with the strongest volume-matched declarative or binary-status control.

\paragraph{H2.} Explicit image-to-object provenance particularly reduces wrong-object and misassociated-photo errors beyond a missing-shot list of comparable length.

\paragraph{H3.} The evidence-rich report does not reduce correct acceptance of complete missions by more than a preregistered non-inferiority margin and improves decision-confidence calibration.

\section{Operationalization and study population}
\subsection{Independent unit of analysis}
One participant-by-mission decision; participant and case are crossed clustering factors. A mission cannot serve as its own displayed ground truth.

\subsection{Data, sensors and ground-truth requirements}
Immutable offline RVI manifests, independent required-shot truth, versioned screenshots and report-generation code, randomized case and display assignments, consent-approved decision logs. Data and models will be versioned with checksums, licensing information, instrument parameters, selection rules and clock provenance. The held-out test cohort must be independent at the acquisition-session, scene or participant level rather than at the individual frame level.

\subsection{Primary and secondary endpoints}
The primary estimand is defined as follows: False-acceptance rate among objectively incomplete mission cases, defined as the count of accepted incomplete cases divided by all evaluated incomplete cases.

The secondary measurements will be: True acceptance of complete cases, wrong-object error identification, decision time, confidence calibration and cognitive workload if justified. The exact numerator, denominator, outcome threshold and permitted exclusions must be fixed in the protocol before inferential analysis.

\section{Planned experiments and research sequence}
\subsection{A. Inspection fixture ground truth}
Create independently adjudicated cases with complete evidence, missing required angles, low-quality shots, wrong-object associations and inconsistent identifiers. Version the requirements and truth labels; balance correct and incorrect missions.

\subsection{B. Interface and sample-size pilot}
Build a browser viewer with response time and confidence capture. Assess comprehension and visual salience, match report verbosity, estimate case/participant variance and select confirmatory sample size by simulated power or required precision.

\subsection{C. Randomized ethics-approved decision experiment}
Randomize report variant and mission case in a balanced incomplete block. Prevent the same participant from seeing the same mission in more than one display format. Collect accept/reject, confidence and latency after required consent.

\subsection{D. Independent confirmatory analysis}
Compare false acceptance on incomplete cases while also reporting true acceptance on complete cases and effects by error type. Analyse a text-length-matched control to rule out verbosity as the only explanatory mechanism.

\section{Comparative design and statistical inference}
Use a mixed-effects logistic model with participant and mission effects, reporting marginal risk differences. Predefine the primary evidence-rich versus strongest report-control contrast, an information-volume control and a non-inferiority margin for complete missions. When human participants or independent trajectories are needed, their number must follow a documented power simulation or precision calculation after a feasible pilot. Report uncertainty, negative results and sensitivity to failure strata. Thresholds, model configurations and exclusion criteria must be frozen before final testing.

\section{Evidence gates, reproducibility and ethical boundary}
Requires prior ethics/consent approvals for human decisions. Scenarios remain browser-based offline fixtures and do not imply Unitree G1 motion or physical inspection results. Each study phase will preserve source SHA or release, configuration, labels, independent oracle, random seeds, raw event timing, and an auditable trace from inputs to computed figures. Simulated evidence is distinct from measured hardware outcomes. Participant studies require prior ethical review and consent.

\section{Results}
Empirical results, tables and uncertainty intervals will be incorporated after execution of the documented evaluation protocol and independent verification of the input evidence.

\section{Discussion}
The planned comparison is intended to determine whether the proposed mechanism yields an identifiable gain over the stated controls. Interpretation must distinguish failure of a scientific hypothesis from an incomplete implementation, and should describe practical operating conditions under which conclusions remain valid.

\section{Conclusion}
This study defines a testable research problem, matched comparators and a primary evaluation target. A substantive empirical conclusion remains conditional on reproducible measurements and appropriate methodological review.

\section{Data and code availability}
The source project \texttt{MatPomGit/arc-robo} remains the authority for the underlying research protocol; this local document is an ARC Robo prospective research plan, not a duplicate final manuscript. Dataset permissions, input manifests, source commits, evaluation scripts and permitted public artifacts will be identified in the final version. No participant records or proprietary instrumentation details are included here.
\begin{thebibliography}{9}
\bibitem{source1} Transparency improves the accuracy of automation use, but automation confidence information does not. Cognitive Research: Principles and Implications (2024). \url{https://doi.org/10.1186/s41235-024-00599-x}.
\bibitem{source2} Trust as a design principle in human-robot collaboration: a review of explainable and adaptive control. Artificial Intelligence Review (2026). \url{https://doi.org/10.1007/s10462-026-11554-1}.
\end{thebibliography}
\end{document}
